% =============================================================================
% CoRaLS Dual-Pol Trigger Signal Chain: Reference Document
% =============================================================================
% Purpose: Detailed technical reference for the simulated signal chain from
%          raw impulse through to S-curve generation. Written for debugging
%          and edification, not publication.
% =============================================================================

\documentclass[11pt,letterpaper]{article}

\usepackage[margin=1.1in]{geometry}
\usepackage{amsmath,amssymb}
\usepackage{fancyvrb}          % Verbatim environments with line numbers
\usepackage{xcolor}
\usepackage{hyperref}
\usepackage{booktabs}
\usepackage{enumitem}
\usepackage{titlesec}
\usepackage{parskip}
\usepackage{mdframed}

% ---- colour scheme for code boxes ----
\definecolor{codebg}{rgb}{0.96,0.96,0.96}
\definecolor{commentcolor}{rgb}{0.4,0.4,0.4}
\definecolor{keywordcolor}{rgb}{0.13,0.29,0.53}
\definecolor{warningbg}{rgb}{1.0,0.95,0.85}
\definecolor{notebg}{rgb}{0.88,0.95,0.88}

\newenvironment{codebox}{%
  \begin{mdframed}[backgroundcolor=codebg,linewidth=0.5pt,innerleftmargin=6pt,innerrightmargin=6pt,innertopmargin=4pt,innerbottommargin=4pt]%
  \small}{\end{mdframed}}

\newenvironment{warningbox}{%
  \begin{mdframed}[backgroundcolor=warningbg,linewidth=1pt,frametitle={\small\textbf{Pitfall / Past Bug}},innerleftmargin=6pt]}%
  {\end{mdframed}}}

\newenvironment{notebox}{%
  \begin{mdframed}[backgroundcolor=notebg,linewidth=0.5pt,frametitle={\small\textbf{Note}},innerleftmargin=6pt]}%
  {\end{mdframed}}}

\RecustomVerbatimEnvironment{Verbatim}{Verbatim}{fontsize=\small,frame=none,numbers=left,numbersep=4pt}

\title{\textbf{CoRaLS Dual-Pol Trigger Simulation}\\\large Signal Chain Reference: From Impulse to S-Curve}
\author{Reference Document --- Not for Publication}
\date{March 2026}

\begin{document}
\maketitle
\tableofcontents
\newpage

% =============================================================================
\section{Overview and Motivation}
% =============================================================================

This document traces the complete simulated signal chain used to generate
efficiency (S-) curves for the CoRaLS dual-polarization radio trigger. The
primary goals are:

\begin{enumerate}[nosep]
  \item Establish a single authoritative reference for how SNR is defined at
        every stage so that comparisons between filter configurations are
        self-consistent.
  \item Document every normalization and threshold convention so that a future
        reader (or a future debugging session) can reproduce results without
        re-reading all source files.
  \item Record bugs that were found and fixed so they are not re-introduced.
\end{enumerate}

The chain consists of the following scripts, roughly in dependency order:

\begin{center}
\begin{tabular}{lll}
\toprule
Script & Purpose & Key output \\
\midrule
\texttt{generate\_dualpol\_noise\_chunked.py} & Generate raw noise files & \texttt{noise/dualpol\_noise\_\{hpol,vpol\}.npy} \\
\texttt{generate\_power\_dualpol.py}           & Beamform noise $\to$ power frames  & \texttt{noise/power\_\{lhcp,rhcp\}\_160\_40*.npy} \\
\texttt{generate\_threshold\_curves\_dualpol.py} & Fit false-rate curve, solve threshold & \texttt{noise/threshold\_analysis\_dualpol*.json} \\
\texttt{snr\_scan\_dualpol.py}                 & Monte Carlo S-curves & \texttt{plots/snr\_scan\_dualpol\_*.npy} \\
\bottomrule
\end{tabular}
\end{center}

Support modules used throughout:
\texttt{payload\_signal.py} (impulse/waveform preparation, beam patterns),
\texttt{coherent\_sum.py} (beamforming, filtering, circular decomposition, power sum),
\texttt{tools/filters.py} (Shannon-Whitaker FIR),
\texttt{noise.py} (thermal noise generator).

% =============================================================================
\section{Noise Generation}
% =============================================================================

\subsection{What is generated}

Noise is generated as band-limited white thermal noise at the RITC ADC sample
rate of $f_s = 4$\,GHz (sample step $\Delta t = 0.25$\,ns) with
$V_\mathrm{rms} = 1.0$.  Eight independent channels are produced: four H-pol
and four V-pol (one per physical antenna location).  The total duration is
$0.5$\,s, giving $2 \times 10^9$ samples per channel.  The arrays are written
directly to \texttt{noise/dualpol\_noise\_hpol.npy} and
\texttt{noise/dualpol\_noise\_vpol.npy} as memory-mapped float64 arrays of
shape \texttt{(4, 2\_000\_000\_000)}.

\begin{codebox}
\begin{Verbatim}
# generate_dualpol_noise_chunked.py  (excerpt)
thermal_noise = noise.ThermalNoise(
    v_rms=vrms,
    fbins=fbins_chunk,
    time_domain_sampling_rate=sample_step_ns   # 0.25 ns
)
# Generated independently for each of 8 channels (4 H + 4 V)
# Saved as memory-mapped float64: shape (4, total_samples)
\end{Verbatim}
\end{codebox}

\begin{notebox}
The raw noise has $V_\mathrm{rms} \approx 1.0$ by construction. After the
first digital lowpass filter (1.5\,GHz Shannon-Whitaker) the RMS drops to
$\sigma_{f1} \approx 0.852$ (measured: 0.8519 H, 0.8524 V). After the
optional second filter at 750\,MHz the RMS drops further to
$\sigma_{f2} \approx 0.592$. These ratios are consistent with
$\sigma_{f1}/\sigma_\mathrm{raw} \approx \sqrt{1.5/2.0} = 0.866$ and
$\sigma_{f2}/\sigma_{f1} \approx \sqrt{0.75/1.5} = 0.707$.
\end{notebox}

% =============================================================================
\section{Impulse Preparation}
% =============================================================================

The template signal is measured impulse data from
\texttt{impulse/corals\_impulse\_sci.txt}.  The key preparation step is the
normalization in \texttt{prepImpulse}:

\begin{codebox}
\begin{Verbatim}
# payload_signal.py -- prepImpulse()
impulse.zeropad(4096)
impulse.fft()
impulse.time = impulse.time - impulse.time[0]

# Normalize: Vpp = 1 exactly
impulse.voltage = impulse.voltage / (
    numpy.max(impulse.voltage) - numpy.min(impulse.voltage)
)
\end{Verbatim}
\end{codebox}

\textbf{Critical consequence:} After \texttt{prepImpulse}, the impulse has
peak-to-peak amplitude $V_{pp} = 1.0$ regardless of any other parameter.
This is the reference against which SNR is defined (see Section~\ref{sec:snr}).

The impulse is not filtered during preparation. The analog (pre-ADC) filter
response is entirely below the Nyquist frequency of the ADC and is
incorporated implicitly in the measured waveform itself.

After applying the Shannon-Whitaker digital filter, the peak-to-peak amplitude
of an individual antenna waveform is essentially unchanged:
$V_{pp}^{f1} \approx V_{pp}^{f2} \approx 1.00$. This is because the impulse's
energy is already well below the filter cutoffs; the filters remove only
out-of-band noise.

% =============================================================================
\section{Per-Antenna Waveform Generation}
% =============================================================================

\texttt{getPayloadWaveforms\_dualpol} generates 8 waveforms (4 H-pol,
4 V-pol) from the single normalized impulse, incorporating:

\begin{itemize}[nosep]
  \item \textbf{Polarization projection.} Signal arrives at polarization angle
        $\psi$ (default 45° for equal power split).
        $\cos\psi$ scales H-pol, $\sin\psi$ scales V-pol.
  \item \textbf{Beam pattern attenuation.}  Each antenna's recorded voltage is
        multiplied by the beam response at the arrival direction
        $(\phi, \theta)$, evaluated separately for E- and H-planes.
  \item \textbf{Geometric delay.} Each antenna's waveform is rolled by
        $\tau_\mathrm{ant}(\phi,\theta)$ samples.
\end{itemize}

\begin{codebox}
\begin{Verbatim}
# payload_signal.py -- getPayloadWaveforms_dualpol() (simplified)
psi_rad = numpy.deg2rad(psi)
pol_h   = numpy.cos(psi_rad)   # H projection coefficient
pol_v   = numpy.sin(psi_rad)   # V projection coefficient

for ant_idx, ant_num in enumerate(antennas):
    delay_samples = int(numpy.round(ant_delay / impulse.dt))

    waveforms[ant_num] = numpy.roll(impulse.voltage * snr, delay_samples) \
        * pol_h \
        * dBtoVoltsAtten(beam_patterns_h[0](phi_interp)) \
        * dBtoVoltsAtten(beam_patterns_h[1](theta_interp))

    waveforms[ant_num + 4] = numpy.roll(impulse.voltage * snr, delay_samples) \
        * pol_v \
        * dBtoVoltsAtten(beam_patterns_v[0](theta_interp)) \
        * dBtoVoltsAtten(beam_patterns_v[1](phi_interp))
\end{Verbatim}
\end{codebox}

The \texttt{snr} argument to \texttt{getPayloadWaveforms\_dualpol} is set to
$1.0$ by \texttt{snr\_scan\_dualpol.py}. Amplitude scaling is applied
externally before signal injection (see Section~\ref{sec:injection}).

At boresight ($\phi = \theta = 0$, $\psi = 45°$), with unit beam response and
zero delay, each antenna delivers $V_{pp} = \cos(45°) = 1/\sqrt{2}$ in H and
$V_{pp} = \sin(45°) = 1/\sqrt{2}$ in V. The quadrature sum gives the full
$V_{pp} = 1$ in power terms.

% =============================================================================
\section{SNR Definition}
\label{sec:snr}
% =============================================================================

\subsection{The definition used in \texttt{snr\_scan\_dualpol.py}}

\begin{mdframed}[backgroundcolor=notebg,linewidth=1pt]
\textbf{Definition:}
\[
  \mathrm{SNR} \;\triangleq\; \frac{V_{pp}^{\mathrm{impulse}}}{2\,\sigma_{f1}}
\]
where $V_{pp}^{\mathrm{impulse}} = 1$ (after \texttt{prepImpulse}) and
$\sigma_{f1}$ is the per-antenna, per-channel noise RMS measured \emph{after}
applying the first 1.5\,GHz Shannon-Whitaker digital filter to the raw noise
files.

This is the \textbf{analog filter reference SNR}: it matches the noise floor
seen by the hardware at the ADC input, after the analog anti-aliasing filter
has acted.
\end{mdframed}

Numerically, $\sigma_{f1} \approx 0.852$ (measured from the noise files), so
at $\mathrm{SNR} = 1$ the signal amplitude scale applied to each antenna is:
\[
  s = \mathrm{SNR} \times 2\,\sigma_{f1} = 1 \times 2 \times 0.852 = 1.704.
\]

\subsection{Why not the raw ($f_s/2 = 2$\,GHz) noise RMS?}

The raw noise RMS $\sigma_\mathrm{raw} \approx 1.0$ (by construction of the
noise generator). It represents the total noise power in the band
$[0, 2\,\mathrm{GHz}]$.  Using $\sigma_\mathrm{raw}$ as the reference would
give $s = 1.0$ at SNR $= 0.5$, introducing a filter-dependent offset on the
x-axis that makes it impossible to compare filt1 vs.\ filt2 directly. Since
the hardware anti-alias filter is always present, $\sigma_{f1}$ is the
operationally meaningful noise floor.

\subsection{Code}

\begin{codebox}
\begin{Verbatim}
# snr_scan_dualpol.py -- noise RMS calculation
rms_sample_size  = 100_000
noise_rms_h = np.sqrt(np.mean(
    noise_h_full[:, :rms_sample_size].astype(np.float64)**2))
noise_rms_v = np.sqrt(np.mean(
    noise_v_full[:, :rms_sample_size].astype(np.float64)**2))
# noise_rms_raw ~ 1.0 by construction

filt_sample_size = 200_000
noise_h_f1 = filters.apply_Shannon_Whitaker_filter(
    noise_h_full[:, :filt_sample_size].astype(np.float64))
noise_v_f1 = filters.apply_Shannon_Whitaker_filter(
    noise_v_full[:, :filt_sample_size].astype(np.float64))
noise_rms_f1_h = np.sqrt(np.mean(noise_h_f1**2))
noise_rms_f1_v = np.sqrt(np.mean(noise_v_f1**2))
noise_rms_f1   = (noise_rms_f1_h + noise_rms_f1_v) / 2
# ~ 0.852

noise_rms = noise_rms_f1  # THIS is the SNR axis denominator
\end{Verbatim}
\end{codebox}

\subsection{Signal injection}
\label{sec:injection}

\begin{codebox}
\begin{Verbatim}
# snr_scan_dualpol.py -- per-trial signal injection
sig_scale     = snr * 2 * noise_rms    # noise_rms = sigma_f1
injected_h    = waveforms_h * sig_scale + event_noise_h
injected_v    = waveforms_v * sig_scale + event_noise_v
\end{Verbatim}
\end{codebox}

At injection SNR $= 1$, each antenna receives:
\[
  v_\mathrm{ant}(t) = 1.704 \cdot w_\mathrm{impulse}(t) + n(t),
  \qquad n \sim \mathcal{N}(0,\,\sigma_\mathrm{raw}^2)
\]
The noise $n(t)$ is the raw (unfiltered) noise segment extracted from the
pre-generated files. It will be filtered during beamforming.

\subsection{Effective SNR after beamforming and filtering}

After 4-antenna coherent beamforming the signal voltage grows by $N_\mathrm{ant}
= 4$, while the incoherent noise power grows by $N_\mathrm{ant}$:
\[
  \mathrm{SNR}_\mathrm{beamformed} = \sqrt{N_\mathrm{ant}} \cdot \mathrm{SNR}_\mathrm{per-antenna}
  = 2 \cdot \mathrm{SNR}.
\]
This improvement applies equally to both filter configurations (since the
beamforming step is the same). The filter only affects the noise floor
\emph{relative to the unfiltered reference}, which is why the SNR axis is
already quoted in $\sigma_{f1}$ units.

% =============================================================================
\section{The Hardware Signal Chain in Software}
\label{sec:chain}
% =============================================================================

\texttt{coherentSum\_dualpol} implements the following ordered steps:

\begin{enumerate}
  \item \textbf{Digitize / decimate} to 4\,GHz ($\Delta t = 0.25$\,ns) by
        integer decimation.
  \item \textbf{Shannon-Whitaker FIR lowpass at 1.5\,GHz} (first filter).
  \item \textbf{Optional Shannon-Whitaker FIR lowpass at 750\,MHz} (second
        filter, enabled by \texttt{apply\_second\_filter=True}).
  \item \textbf{Coherent sum} each polarization:
        apply geometric delays (integer-sample roll) and sum all antennas.
  \item \textbf{Circular polarization conversion:}
        $\mathrm{LHCP} = (H + jV)/\sqrt{2}$;
        $\mathrm{RHCP} = (H - jV)/\sqrt{2}$, done in the Fourier domain.
\end{enumerate}

\begin{codebox}
\begin{Verbatim}
# coherent_sum.py -- coherentSum_dualpol() (key steps)

# Step 1: Decimate to 4 GHz
dt_in          = timebase[1] - timebase[0]
decimate_factor = int(aso_geometry.ritc_sample_step / dt_in)
waveforms_h_dig = waveforms_h[:, ::decimate_factor]
waveforms_v_dig = waveforms_v[:, ::decimate_factor]

# Step 2: First filter (1.5 GHz Shannon-Whitaker)
waveforms_h_filt1 = filters.apply_Shannon_Whitaker_filter(waveforms_h_dig)
waveforms_v_filt1 = filters.apply_Shannon_Whitaker_filter(waveforms_v_dig)

# Step 2b: Optional second filter (750 MHz)
if apply_second_filter:
    waveforms_h_filt2 = filters.apply_Shannon_Whitaker_filter(
        waveforms_h_filt1, fc=fc_second)
    waveforms_v_filt2 = filters.apply_Shannon_Whitaker_filter(
        waveforms_v_filt1, fc=fc_second)
else:
    waveforms_h_filt2 = waveforms_h_filt1
    waveforms_v_filt2 = waveforms_v_filt1

# Step 3: Coherent sum
coh_h, tb = coherentSum(waveforms_h_filt2, timebase, delays, downsample=False)
coh_v, _  = coherentSum(waveforms_v_filt2, timebase, delays, downsample=False)

# coherentSum() applies integer-sample roll per antenna:
#   _delay = -int(round(delays[ant] / dt))
#   coh_sum += roll(waveforms[ant], _delay)

# Step 4-6: Circular basis conversion via FFT
H_fft    = numpy.fft.rfft(coh_h)
V_fft    = numpy.fft.rfft(coh_v)
V_shifted = 1j * V_fft          # pi/2 phase shift
LHCP_fft  = (H_fft + V_shifted) / numpy.sqrt(2)
RHCP_fft  = (H_fft - V_shifted) / numpy.sqrt(2)
lhcp = numpy.fft.irfft(LHCP_fft, n=len(coh_h))
rhcp = numpy.fft.irfft(RHCP_fft, n=len(coh_h))
\end{Verbatim}
\end{codebox}

\noindent\textbf{Order matters.} The filter is applied \emph{per antenna
before beamforming}. This matches the hardware implementation where each
channel is filtered digitally before the combinatorial sum. Apply the filter
after the sum and you will get a different (incorrect) threshold.

% =============================================================================
\section{Sliding-Window Power Sum}
% =============================================================================

\begin{codebox}
\begin{Verbatim}
# coherent_sum.py -- powerSum()
def powerSum(coh_sum, window=32, step=16):
    num_frames      = int((len(coh_sum) - window) // step) + 1
    coh_power       = numpy.abs(coh_sum)**2
    coh_sum_windowed = numpy.lib.stride_tricks.as_strided(
        coh_power,
        shape   = (num_frames, window),
        strides = (coh_power.strides[0] * step, coh_power.strides[0])
    )
    power = numpy.sum(coh_sum_windowed, axis=1) / window
    return power.astype(numpy.float64), num_frames
\end{Verbatim}
\end{codebox}

For the trigger, the default parameters are:
\[
  W = 160\;\mathrm{samples} = 40\;\mathrm{ns},
  \qquad
  S = 40\;\mathrm{samples} = 10\;\mathrm{ns}.
\]
The returned quantity is the mean squared voltage per window:
\[
  P[k] = \frac{1}{W} \sum_{i=kS}^{kS+W-1} \bigl|v[i]\bigr|^2.
\]
This is \emph{not} normalized by $N_\mathrm{ant}$.  The beamformed voltage
at boresight is $\sum_n v_n$, so for pure signal power grows as
$N_\mathrm{ant}^2$ and for pure noise as $N_\mathrm{ant}$ under coherent and
incoherent sums, respectively.

\begin{warningbox}
\textbf{Do not divide by $N_\mathrm{ant}$ before comparing against the
threshold.}  The threshold values in the JSON files (e.g.\ $T_{f1} = 5.006$
at 0.1\,Hz) are calibrated on the un-normalized beamformed noise power using
\texttt{generate\_power\_dualpol.py}, which never divides by
$N_\mathrm{ant}$. Dividing the scan power by $N_\mathrm{ant}$ before the
comparison shrinks the effective threshold by a factor of $N_\mathrm{ant}$,
making it $4\times$ too strict (old bug, fixed March 2026).
\end{warningbox}

% =============================================================================
\section{Threshold Calibration}
% =============================================================================

\subsection{Noise power distribution}

\texttt{generate\_power\_dualpol.py} runs the full beamforming chain on the
pre-generated noise files (0.5\,s) and outputs per-frame power values for
LHCP and RHCP.  At boresight $(\phi=0°, \theta=-30°)$ with $N_\mathrm{ant}=4$
and no signal, the mean power per frame for the first-filter configuration is:
\[
  \langle P_\mathrm{noise}^{f1} \rangle
    = N_\mathrm{ant} \cdot \sigma_{f1}^2
    = 4 \times 0.852^2
    \approx 2.90.
\]
For the second-filter configuration:
\[
  \langle P_\mathrm{noise}^{f2} \rangle
    = N_\mathrm{ant} \cdot \sigma_{f2}^2
    = 4 \times 0.592^2
    \approx 1.40.
\]
This is consistent with the debug print during the SNR scan showing
$\mathrm{pow} \approx 3.3$ and $\mathrm{pow} \approx 1.7$ at low SNR
(including the signal term even at SNR $= 0.2$).

\subsection{False-rate curve fit}

\texttt{generate\_threshold\_curves\_dualpol.py} scans threshold $T$ and
counts the fraction of frames exceeding $T$, converting to a rate using the
frame step time $\Delta t_\mathrm{frame} = S / f_s = 10$\,ns.  The dependence
is fit to:
\[
  R(T) = A\,e^{-BT}.
\]

\subsection{Global rate constraint and threshold solving}

The global coincidence rate with $N_\mathrm{beams}$ independent beams is:
\[
  R_\mathrm{global} = N_\mathrm{beams} \cdot R_\mathrm{coinc/beam},
  \qquad
  R_\mathrm{coinc/beam} = 2\,R_L(T)\,R_R(T)\,\tau,
\]
where $\tau = 20$\,ns is the coincidence window (set at threshold generation
time) and the factor 2 comes from the Poisson approximation for two
independent processes.  For $N_\mathrm{beams} = 63$ (from
\texttt{plots/optimal\_beams\_sprint6.npy}) and a target of 0.1\,Hz:

\begin{codebox}
\begin{Verbatim}
# generate_threshold_curves_dualpol.py
tau_s         = args.coincidence_window * 1e-9    # 20 ns -> 2e-8 s
target_per_beam = target / args.n_beams           # 0.1 / 63 = 1.587e-3 Hz
# Symmetric solution: T_L = T_R = T
# 2 * A^2 * tau * exp(-2*B*T) = target_per_beam
# => T = -ln(target_per_beam / (2*A_L*A_R*tau)) / (B_L + B_R)
arg  = target_per_beam / (2 * A_lhcp * A_rhcp * tau_s)
T_sym = -numpy.log(arg) / (B_lhcp + B_rhcp)
\end{Verbatim}
\end{codebox}

Results for $0.1$\,Hz global rate:

\begin{center}
\begin{tabular}{lcc}
\toprule
Configuration & $T$ (power units) & $\sigma_{f1}$ referred \\
\midrule
filt1 (1.5\,GHz) & 5.006 & $\approx 1.31\,\sigma_\mathrm{beam}$ \\
filt2 (750\,MHz) & 2.981 & $\approx 1.31\,\sigma_\mathrm{beam}$ \\
\bottomrule
\end{tabular}
\end{center}

The thresholds are different in absolute units but represent the same
statistical operating point because the noise floor they are measured against
is different ($\sigma_{f2} < \sigma_{f1}$).

% =============================================================================
\section{S-Curve Generation}
% =============================================================================

\subsection{Monte Carlo loop}

\texttt{snr\_scan\_dualpol.run\_snr\_scan\_dualpol} iterates over an SNR grid
and, for each point, runs $N_\mathrm{trials} = 1000$ independent signal+noise
trials. For each trial:

\begin{enumerate}
  \item Extract a random $N_\mathrm{samp}$-sample noise segment (same start
        index for all 4 antennas in that trial, independent draw per trial).
  \item Scale the signal template and add noise:
        $v_\mathrm{ant}^{H,V}(t) = s \cdot w_\mathrm{ant}^{H,V}(t) + n_\mathrm{ant}^{H,V}(t)$.
  \item Run \texttt{coherentSum\_dualpol} $\to$ \texttt{powerSum} $\to$
        compare against threshold.
  \item Record LHCP hit, RHCP hit, and coincidence hit.
\end{enumerate}

\begin{codebox}
\begin{Verbatim}
# snr_scan_dualpol.py -- inner trial loop
sig_scale   = snr * 2 * noise_rms        # noise_rms = sigma_f1
injected_h  = waveforms_h * sig_scale + event_noise_h
injected_v  = waveforms_v * sig_scale + event_noise_v

lhcp, rhcp, _ = trigger.coherentSum_dualpol(
    injected_h, injected_v, timebase, delays_array,
    downsample=True, apply_filter=True, digitize_first=True,
    apply_second_filter=apply_second_filter,
    fc_second=payload.FC_SECOND_FILTER,
    output='circular'
)

power_lhcp, _ = trigger.powerSum(lhcp, window=window, step=step)
power_rhcp, _ = trigger.powerSum(rhcp, window=window, step=step)
# NO division by n_ant here

max_lhcp = np.max(power_lhcp)
max_rhcp = np.max(power_rhcp)

lhcp_triggered = max_lhcp > threshold_lhcp
rhcp_triggered = max_rhcp > threshold_rhcp

if lhcp_triggered:  hits_lhcp += 1
if rhcp_triggered:  hits_rhcp += 1

# Coincidence: frame-by-frame AND gate
if lhcp_triggered and rhcp_triggered:
    coinc_frames = (power_lhcp > threshold_lhcp) & (power_rhcp > threshold_rhcp)
    if np.any(coinc_frames):
        hits_coinc += 1
\end{Verbatim}
\end{codebox}

\subsection{Trigger logic}

\subsubsection{Single-polarization trigger}

An event is flagged as LHCP-triggered if the maximum power over all frames
exceeds the threshold:
\[
  \mathtt{lhcp\_triggered} = \max_k P_\mathrm{LHCP}[k] > T_{f1}.
\]

\subsubsection{Coincidence trigger}

\textbf{Current implementation (frame-by-frame AND):} An event passes
coincidence if there exists at least one frame $k$ where both channels
simultaneously exceed threshold:
\[
  \mathtt{coinc} = \exists\,k:\;
  P_\mathrm{LHCP}[k] > T \;\wedge\; P_\mathrm{RHCP}[k] > T.
\]

\begin{warningbox}
\textbf{Do not use argmax-based coincidence for LHCP/RHCP.}
An earlier implementation found the peak frame index in each channel
independently and checked whether those frame indices differed by less than
$\tau_\mathrm{coinc}$ in time.  This is incorrect here because LHCP and RHCP
are derived from the same beamformed H-pol and V-pol streams --- there is no
physical time offset between them.  The argmax can differ between LHCP and
RHCP due to noise, especially when the 750\,MHz filter disperses the impulse
response over more frames (the peak is lower and broader), causing the per-channel
noise fluctuations to shift the argmax index away from the true signal peak.
This produced a spurious coincidence efficiency plateau at $\sim$50--70\% at
high SNR for the second-filter configuration (old bug, fixed March 2026).
\end{warningbox}

\subsubsection{Why LHCP and RHCP are derived from the same array}

The circular decomposition is done on the coherently beamformed H and V
streams, which come from the same antenna array observing the same sky
direction in the same time window.  There is therefore no geometrically
induced time offset between LHCP and RHCP; any apparent timing difference is
purely due to noise.

% =============================================================================
\section{Filter Comparison}
% =============================================================================

\subsection{What the second filter does}

The 750\,MHz lowpass is applied per-antenna after the 1.5\,GHz first filter.
Its effect on the impulse response is negligible ($V_{pp}$ is unchanged; all
impulse energy is already below 750\,MHz), but the noise RMS drops by a
further factor of $\approx 0.695$ (ratio $\sigma_{f2}/\sigma_{f1}$,
consistent with $\sqrt{750/1500} = 0.707$).

This means that at the same per-antenna SNR (defined via $\sigma_{f1}$), the
signal-to-noise ratio in the power domain improves by:
\[
  \frac{(V_{pp}^{f2}/\sigma_{f2})^2}{(V_{pp}^{f1}/\sigma_{f1})^2}
  = \left(\frac{\sigma_{f1}}{\sigma_{f2}}\right)^2
  \approx \left(\frac{0.852}{0.592}\right)^2
  \approx 2.07.
\]
This\textasciitilde{}2$\times$ improvement in power SNR corresponds to about
$\sqrt{2} \approx 1.44\times$ lower SNR required for a given detection
probability.

\subsection{Why this improvement does not show up as an x-axis shift on the S-curve}

Because the SNR axis is defined in $\sigma_{f1}$ units, the filter improvement
is fully captured in the threshold value (2.981 vs.\ 5.006).  Therefore the
S-curve for filt2 occupies approximately the same x-axis range as filt1; the
\emph{performance improvement is expressed as a lower threshold at the same
false rate}, not as a leftward shift of the S-curve.

\subsection{How to present both curves with a visible leftward shift}

To show the filt2 improvement as a leftward shift on the same axes, rescale
the filt2 x-axis values so that the denominator changes from $\sigma_{f1}$ to
$\sigma_{f2}$:

\[
  \mathrm{SNR}_{f2\text{-axis}} = \mathrm{SNR}_{f1\text{-axis}} \times \frac{\sigma_{f2}}{\sigma_{f1}}
  \approx \mathrm{SNR}_{f1\text{-axis}} \times 0.695.
\]

\noindent Conceptually: the filt2 scan injects signal at amplitude
$s = \mathrm{SNR} \times 2\sigma_{f1}$, but the filt2 system's own noise floor
is $\sigma_{f2}$. So the \emph{effective} per-antenna SNR as seen by the filt2
system is:
\[
  \frac{s}{2\sigma_{f2}} = \mathrm{SNR}_{f1\text{-axis}} \times \frac{\sigma_{f1}}{\sigma_{f2}}
  \approx \mathrm{SNR}_{f1\text{-axis}} \times 1.44.
\]
Re-expressing the filt2 x-axis in these native units compresses it by
$\sigma_{f2}/\sigma_{f1} \approx 0.695$, shifting the curve leftward by
$\approx 1.44\times$.

In Python, given \texttt{snr\_arr} is the common SNR array from the scan:

\begin{codebox}
\begin{Verbatim}
sigma_f1 = 0.852  # measured
sigma_f2 = 0.592  # measured

# filt1 curve: x-axis unchanged
ax.plot(snr_arr, eff_coinc_f1, label='filt1 (1.5 GHz)')

# filt2 curve: rescale x-axis to filt2 native SNR units
snr_arr_f2 = snr_arr * (sigma_f2 / sigma_f1)   # ~ 0.695 * snr_arr
ax.plot(snr_arr_f2, eff_coinc_f2, label='filt2 (750 MHz, sigma_f2 axis)')

ax.set_xlabel(r'per-antenna $V_{pp} / (2\sigma)$ [own filter reference]')
\end{Verbatim}
\end{codebox}

The two curves will then cross the 50\,\% line at different x positions
($\approx 1.25$ for filt1 vs.\ $\approx 0.9$ for filt2), making the gain from
the second filter visually apparent.  The trade-off is that the two curves now
have \emph{different} x-axis denominators, so the plot does not directly answer
``at what injected amplitude does each system trigger?'' --- it answers ``how
many multiples of each system's own noise floor are needed?''

\begin{notebox}
The choice of presentation depends on the audience.  A \emph{common} $\sigma_{f1}$
axis (what the scan currently produces) is most useful when comparing with a
fixed injected signal level (e.g.\ a real event of known voltage). The
\emph{rescaled} presentation is most useful for comparing the intrinsic
sensitivity of two filter configurations on equal footing.  Both are physically
correct.
\end{notebox}

% =============================================================================
\section{Known Bugs and Fixes (History)}
% =============================================================================

\begin{description}[style=nextline]

\item[Power normalized by $N_\mathrm{ant}$ (March 2026 fix)]
The SNR scan divided \texttt{power\_lhcp} and \texttt{power\_rhcp} by
$N_\mathrm{ant} = 4$ before comparing to the threshold.  The threshold JSON
values are calibrated on un-normalized beamformed noise power (from
\texttt{generate\_power\_dualpol.py}, which never normalizes).  Dividing the
scan power by 4 made the effective threshold 4$\times$ too stringent,
requiring about $4\times$ more signal power (i.e.\ $2\times$ more SNR on the
voltage axis) to trigger. This inflated the 50\,\% efficiency point from
$\approx$\,SNR\,1.25 to $\approx$\,SNR\,5--6. The comment in the code claimed
the normalization was needed to express signal relative to a single-antenna
threshold, which was wrong: the thresholds are already set for a
4-antenna beamformed system.

\item[SNR reference used raw noise RMS $\sigma_\mathrm{raw} \approx 1.0$ (March 2026 fix)]
Early versions used the full-band noise RMS as the injection reference.  Since
$\sigma_\mathrm{raw} \approx 1.0$ and $V_{pp} = 1$, this accidentally produced
$s \approx 2 \times 1 = 2.0$ at SNR $= 1$, injecting more power than intended
relative to the filter-limited noise floor.  Changing to $\sigma_{f1}$ gives a
physically meaningful SNR that corresponds to what the hardware ADC samples
after the analog anti-aliasing filter.

\item[Argmax-based coincidence window for LHCP/RHCP (March 2026 fix)]
The coincidence check found the time of the peak power frame independently for
LHCP and RHCP and required them to agree within $\tau_\mathrm{coinc}$.  With
the 750\,MHz second filter, the impulse is dispersed over $\sim$5--7 frames
with a flatter peak; noise fluctuations then shift the argmax within the
trigger window, causing LHCP and RHCP to report different peak times even
though both channels were triggered by the same physical signal.  This produced
a coincidence efficiency plateau at $\sim$50--70\,\% at high SNR for the
second-filter configuration. The fix is to use a frame-by-frame AND gate,
which checks whether any single frame simultaneously exceeds threshold in both
channels.

\item[Hard-coded \texttt{xlim=[0,100]} in \texttt{payload\_signal.py} panel 10d]
The power-sum step plot was displayed with a fixed x-axis of 0--100\,ns.  The
actual number of visible power frames (each 10\,ns wide) was different, so the
annotation counts for threshold-crossing frames did not match the visible
display. Fixed by computing the x-axis extent dynamically.

\item[Step plot missing last frame (\texttt{where='post'})]
A step plot with $N$ frame timestamps and \texttt{where='post'} only renders
$N-1$ visible bars (the last bar has no ``next'' timestamp to define its right
edge).  Fixed by appending a closing timestamp at
$t_\mathrm{last} + \Delta t_\mathrm{step}$.

\item[n\_beams defaulted to 100]
The threshold generation script defaulted to $N_\mathrm{beams} = 100$ as a
placeholder. The actual number from the beam optimization is 63 (from
\texttt{plots/optimal\_beams\_sprint6.npy}).  Added \texttt{--beam-data-file}
argument that reads this value automatically.

\end{description}

% =============================================================================
\section{Numerical Summary of Key Constants}
\label{sec:constants}
% =============================================================================

\begin{center}
\begin{tabular}{lll}
\toprule
Quantity & Value & Source \\
\midrule
ADC sample rate $f_s$ & 4\,GHz & \texttt{tools/CoRaLs\_geometry.py} \\
ADC sample step $\Delta t$ & 0.25\,ns & \\
First filter cutoff $f_{c,1}$ & 1.5\,GHz & Shannon-Whitaker FIR \\
Second filter cutoff $f_{c,2}$ & 750\,MHz & Shannon-Whitaker FIR \\
Power window $W$ & 160 samples = 40\,ns & \\
Power step $S$   & 40 samples = 10\,ns  & \\
Coincidence window $\tau_\mathrm{coinc}$ (threshold gen.) & 20\,ns & \\
$N_\mathrm{beams}$ & 63 & \texttt{plots/optimal\_beams\_sprint6.npy} \\
Target global rate & 0.1\,Hz & \\
\midrule
$\sigma_\mathrm{raw}$ (per-antenna, pre-filter) & $\approx 0.999$ & measured \\
$\sigma_{f1}$ (per-antenna, post 1.5\,GHz filter) & $\approx 0.852$ & measured \\
$\sigma_{f2}$ (per-antenna, post 750\,MHz filter)  & $\approx 0.592$ & measured \\
$V_{pp}$ after \texttt{prepImpulse} & 1.0 (exact) & normalized \\
$V_{pp}$ after filt1 per antenna & $\approx 1.000$ & measured \\
$V_{pp}$ after filt2 per antenna & $\approx 1.000$ & measured \\
\midrule
Threshold $T_{f1}$ (filt1, 0.1\,Hz) & 5.006 & fitted \\
Threshold $T_{f2}$ (filt2, 0.1\,Hz) & 2.981 & fitted \\
\midrule
SNR\,50\% coinc efficiency (filt1) & $\approx 1.25$ & scan \\
SNR\,50\% coinc efficiency (filt2) & $\approx 1.30$ & scan \\
SNR\,100\% coinc efficiency (both) & $\approx 1.8$ & scan \\
\bottomrule
\end{tabular}
\end{center}

% =============================================================================
\section{Running the Full Chain}
% =============================================================================

\begin{codebox}
\begin{Verbatim}
# 1. Generate noise (run once; ~0.5 s of data, ~16 GB on disk)
python generate_dualpol_noise_chunked.py --duration 0.5 --chunk-duration 0.01

# 2. Generate beamformed noise power files (filt1 only)
python generate_power_dualpol.py --window 160 --step 40

# 2b. And with second filter
python generate_power_dualpol.py --window 160 --step 40 \
    --apply-second-filter --suffix "_filt2_750MHz"

# 3. Calibrate thresholds (filt1)
python generate_threshold_curves_dualpol.py \
    --beam-data-file plots/optimal_beams_sprint6.npy \
    --window 160 --step 40 \
    --target-rate 0.1 1.0

# 3b. Calibrate thresholds (filt2)
python generate_threshold_curves_dualpol.py \
    --beam-data-file plots/optimal_beams_sprint6.npy \
    --window 160 --step 40 \
    --suffix "_filt2_750MHz" \
    --target-rate 0.1 1.0

# 4. Run S-curve comparison
python snr_scan_dualpol.py \
    --filter-scan \
    --threshold-json noise/threshold_analysis_dualpol.json \
    --filt2-threshold-json noise/threshold_analysis_dualpol_filt2_750MHz.json \
    --threshold-index 0
\end{Verbatim}
\end{codebox}

\end{document}
